
\begin{exercicio}
  Resolva as equações abaixo.
  \begin{mathlist}[1]
        \item 4 \left(x - \frac{3}{4}\right) (x - 6) = 0
        \item (x - 9)^2 = 0
        \item (x - 5) (2 - x) = 0
        \item 4x (x + 8) = 0
        \item 8 \left(x + \frac{1}{2}\right) (x - 4) = 0
        \item (5x + 3) (2x + 7) = 0
        \item \left(\frac{x}{4} - \frac{3}{2}\right) \left(3 - \frac{x}{4}\right) = 0
        \item \sqrt{2} \left(x + \sqrt{2}\right) \left(\frac{x}{\sqrt{2}} - 1\right) = 0
  \end{mathlist}
  \begin{answers}
      \begin{mathlist}[2]
        \item x_1 = \sfrac{3}{4} \qquad x_2 = 6
        \item x_1 = 9
        \item x_1 = 2 \qquad x_2 = 5
        \item x_1 = -8 \qquad x_2 = 0
        \item x_1 = - \sfrac{1}{2} \qquad x_2 = 4
        \item x_1 = - \sfrac{7}{2} \qquad x_2 = - \sfrac{3}{5}
        \item x_1 = 6 \qquad x_2 = 12
        \item x_1 = \sqrt{2} \qquad x_2 = - \sqrt{2}
    \end{mathlist}
  \end{answers}
\end{exercicio}

\begin{exercicio}
    Reescreva as equações do exercício anterior na forma $ax^2 + bx + c = 0$.
    \begin{answers}
        \begin{mathlist}[2]
            \item 4 x^{2} - 27 x + 18 = 0 
            \item x^{2} - 18 x + 81 = 0 
            \item - x^{2} + 7 x - 10 = 0 
            \item 4 x^{2} + 32 x = 0 
            \item 8 x^{2} - 28 x - 16 = 0 
            \item 10 x^{2} + 41 x + 21 = 0 
            \item - \frac{x^{2}}{16} + \frac{9 x}{8} - \frac{9}{2} = 0 
            \item x^{2} - 2 = 0 
        \end{mathlist}
    \end{answers}
\end{exercicio}

\begin{exercicio}
    Encontre as soluções reais das equações abaixo, caso existam.
    \begin{mathlist}[2]
        \item x^2 - 10 = 0
        \item 3x^2 - 75 = 0
        \item 4x^2 + 81 = 0
        \item \frac{x^2}{6} - \frac{24}{9} = 0
        \item (x - 2)^2 = 4^2
        \item (2x - 1)^2 - 25 = 0
        \item (x + 3)^2 = \frac{1}{9}
        \item \left(\frac{x}{2} + 1\right)^2 = \frac{9}{4}
    \end{mathlist}
    \begin{answers}
        \begin{mathlist}[2]
            \item x_1 = \sqrt{10} \quad x_2 =\!-\!\sqrt{10}
            \item x_1 = -5 \qquad x_2 = 5 
            \item \) Não há soluções
            \item x_1 = -4 \qquad x_2 = 4 
            \item x_1 = -2 \qquad x_2 = 6 
            \item x_1 = -2 \qquad x_2 = 3 
            \item x_1 = - \sfrac{10}{3} \quad x_2 = -\sfrac{8}{3}
            \item x_1 = -5 \qquad x_2 = 1 
        \end{mathlist}
    \end{answers}
\end{exercicio}

\begin{exercicio}
    Determine as raízes das equações.
    \begin{mathlist}[2]
        \item x^2 - 4x = 0
        \item 5x^2 + x = 0
        \item x^2 = -7x
        \item 2x^2 - 3x = 0
        \item -3x^2 - \frac{x}{2} = 0
        \item \frac{x^2}{3} - \frac{x}{6} = 0
        \item 2x - \sqrt{2}x^2 = 0
        \item \sqrt{3}x - \frac{x^2}{\sqrt{3}} = 0
    \end{mathlist}
    \begin{answers}
        \begin{mathlist}[2]
            \item x_1 = 0 \qquad x_2 = 4 
            \item x_1 = - \sfrac{1}{5} \qquad x_2 = 0 
            \item x_1 = -7 \qquad x_2 = 0 
            \item x_1 = 0 \qquad x_2 = \sfrac{3}{2} 
            \item x_1 = - \sfrac{1}{6} \qquad x_2 = 0 
            \item x_1 = 0 \qquad x_2 = \sfrac{1}{2} 
            \item x_1 = 0 \qquad x_2 = \sqrt{2} 
            \item x_1 = 0 \qquad x_2 = 3 
        \end{mathlist}
    \end{answers}
\end{exercicio}

\begin{exercicio}
    Usando a fórmula de Bhaskara, determine, quando possível, as raízes reais das equações.
    \begin{mathlist}[2]
        \item x^2 - 6x + 8 = 0
        \item x^2 - 2x - 15 = 0
        \item x^2 + 6x + 9 = 0
        \item x^2 + 8x + 12 = 0
        \item 2x^2 + 8x - 10 = 0
        \item x^2 - 6x + 10 = 0
        \item 2x^2 - 7x - 4 = 0
        \item 6x^2 - 5x + 1 = 0 
        \item x^2 - 4x + 13 = 0
        \item 25x^2 - 20x + 4 = 0
        \item x^2 - 2\sqrt{5}x + 5 = 0
        \item 2x^2 - 2\sqrt{2}x - 24 = 0
        \item 3x^2 -0,3x - 0,36 = 0
        \item x^2 - 2,4x + 1,44 = 0
        \item x^2 + 2x + 5 = 0
        \item(x + 8)^2 + 4x = 0
        \item(3 - 4x) (x + 3) = 9
        \item (2 - 3x) (2x - 5) = 4
    \end{mathlist}    
    \begin{answers}
        \begin{mathlist}[2]
            \item x_1 = 2 \qquad x_2 = 4 
            \item x_1 = -3 \qquad x_2 = 5 
            \item x_1 = -3 
            \item x_1 = -6 \qquad x_2 = -2 
            \item x_1 = -5 \qquad x_2 = 1 
            \item \) Não há soluções
            \item x_1 = - \sfrac{1}{2} \qquad x_2 = 4 
            \item x_1 = \sfrac{1}{3} \qquad x_2 = \sfrac{1}{2} 
            \item \) Não há soluções
            \item x_1 = \sfrac{2}{5} 
            \item x_1 = \sqrt{5} 
            \item x_1 = - 2 \sqrt{2} \qquad x_2 = 3 \sqrt{2} 
            \item x_1 = -0,3 \qquad x_2 = 0,4 
            \item \) Não há soluções
            \item \) Não há soluções
            \item x_1 = -16 \qquad x_2 = -4 
            \item x_1 = - \sfrac{9}{4} \qquad x_2 = 0 
            \item x_1 = \sfrac{7}{6} \qquad x_2 = 2 
        \end{mathlist}
    \end{answers}
\end{exercicio}

\begin{exercicio}
    Determine quantas raízes as equações abaixo possuem.
    \begin{mathlist}[2]
        \item 2x^2 + 12x + 18 = 0
        \item x^2 - 3x + 8 = 0
        \item -2x^2 - 5x + 9 = 0
        \item \frac{x^2}{5} - 2x + 20 = 0
        \item -x^2 + 16x - 64 = 0
        \item 3x^2 - 4x + 1 = 0
    \end{mathlist}      
    \begin{answers}
        \begin{mathlist}[2]
            \item n=1 \qquad x_1 = -3 
            \item n=0 
            \item n=2 \quad  x_1 = \frac{-5 \pm\sqrt{97}}{4}
            \item n=0 
            \item n=1 \qquad x_1 = 8 
            \item n=2 \qquad x_1 = \sfrac{1}{3} \quad x_2 = 1 
        \end{mathlist}
    \end{answers}
\end{exercicio}

